\makeatletter
\def\input@path{{./}{../}{../../}{../../../}{../../../../}{preamble/}{../preamble/}{../../preamble/}{../../../preamble/}}
\makeatother
\input{preamble_exams.tex}
\input{graphs/highlights.tex}
% \GTenHighlightsfalse % Uncomment to disable instructional highlights.
\input{graphs/specific_highlights.tex}

\hypersetup{pdftitle={Grade 10 Global Economics: Term 3, Lesson 1, C6},pdfsubject={Home practice: Standardisation in Economic Models}}
\title{Grade 10 Global Economics\\Term 3, Lesson 1: Home Practice\\C6 --- Standardisation in Economic Models}
\author{}
\date{}
\begin{document}
\maketitle
\setcounter{tocdepth}{1}
\setcounter{secnumdepth}{2}
\phantomsection\label{toc}
\tableofcontents
\section*{Evaluated criteria}
\begin{description}
\item[\textbf{C6}] Uses standardisation \(z=(x-\mu)/\sigma\) to transform normal random variables and solve probability problems involving real-life economic and financial data.
\end{description}
\textbf{How to practise.} Write the mean and standard deviation with units first. Take the square root if variance is supplied, and standardise every boundary separately. Try each task on paper before opening its solution link. Record the model, working, units and an interpretation; use the arrows to return.
All datasets and economic or financial models in this activity are illustrative, rather than observed market evidence.

\clearpage
\section{Introduction}\label{sec:introduction}
\subsection*{Highlight legend}
\phantomsection\label{tab:highlight-legend}
\begin{tabularx}{\linewidth}{@{}>{\raggedright\arraybackslash}p{.20\linewidth}>{\raggedright\arraybackslash}p{.22\linewidth}>{\raggedright\arraybackslash}X@{}}
\toprule
\textbf{Symbol or notation} & \textbf{Concept} & \textbf{Brief meaning}\\
\midrule
$\GTenHLValue{X}\sim N(\mu,\sigma^2)$ & \GTenHLValue{Normal variable} & Model in original units. \\
$\HLPopulationMean{\mu}$ & \HLPopulationMean{Mean} & Population centre. \\
$\HLPopulationSD{\sigma},\ \HLPopulationVariance{\sigma^2}$ & \HLPopulationSD{Standard deviation} and \HLPopulationVariance{variance} & Spread and squared spread. \\
$\HLCriticalValue{z}=(x-\mu)/\sigma$ & \HLCriticalValue{Standard score} & Signed distance in SD units. \\
$\HLCumulativeProbability{\Phi(z)}$ & \HLCumulativeProbability{Left area} & Standard-normal cumulative probability. \\
\bottomrule
\end{tabularx}
\subsection{Introductory example 1: Weekly household spending}\label{ex:intro-one}
Illustrative weekly spending $X\sim N(100,20^2)$ dollars has \HLPopulationMean{mean} $\HLPopulationMean{\mu=100}$ dollars and \HLPopulationSD{standard deviation} $\HLPopulationSD{\sigma=20}$ dollars. A value $x=120$ corresponds to $\HLCriticalValue{z=(120-100)/20=1}$: one standard deviation above the mean.
\subsection{Introductory example 2: Variance is not the divisor}\label{ex:intro-two}
An illustrative monthly return $R\sim N(2,3^2)$ is measured in percentage points. Its \HLPopulationVariance{variance} is $\HLPopulationVariance{9}$ squared percentage points, while its standard deviation is $3$ percentage points. A zero return has $z=(0-2)/3=-2/3$, below its expected return. Dividing by nine would use the wrong quantity and units.
\section{Notation Reference and Practice Procedure}\label{sec:notation-and-procedure}
\subsection{Notation reference}\label{tab:notation-reference}
\begin{tabularx}{\linewidth}{@{}>{\raggedright\arraybackslash}p{.16\linewidth}>{\raggedright\arraybackslash}p{.18\linewidth}>{\raggedright\arraybackslash}p{.27\linewidth}>{\raggedright\arraybackslash}X@{}}
\toprule
\textbf{Symbol or notation} & \textbf{Concept} & \textbf{Definition} & \textbf{Meaning in the introductory examples}\\
\midrule
$\GTenHLValue{X}\sim N(\mu,\sigma^2)$ & \GTenHLValue{Normal variable} & Here the second normal parameter is variance; the model is continuous. & \textbullet\ \hyperref[ex:intro-one]{Example 1}: Weekly spending is in dollars. \\
$\HLPopulationMean{\mu}$ & \HLPopulationMean{Mean} & Model average in the same units as $X$. & \textbullet\ \hyperref[ex:intro-one]{Example 1}: Spending is centred at 100 dollars. \\
$\HLPopulationSD{\sigma},\ \HLPopulationVariance{\sigma^2}$ & \HLPopulationSD{Standard deviation} and \HLPopulationVariance{variance} & $sigma=sqrt{sigma^2}>0$ has original units; variance has squared units. & \textbullet\ \hyperref[ex:intro-two]{Example 2}: Return spread is 3 percentage points; variance is 9 squared percentage points. \\
$\HLCriticalValue{z}=(x-\mu)/\sigma$ & \HLCriticalValue{Standard score} & Subtract the mean, then divide by standard deviation; units cancel. & \textbullet\ \hyperref[ex:intro-one]{Example 1}: 120 dollars is one SD above the spending mean. \\
$\HLCumulativeProbability{\Phi(z)}$ & \HLCumulativeProbability{Left area} & For $Zsim N(0,1)$, $P(Zle z)$. & \textbullet\ \hyperref[ex:intro-one]{Example 1}: Used after converting a spending boundary. \\
\bottomrule
\end{tabularx}
\subsection{Practice procedure}
\begin{enumerate}
\item Choose the stated normal model and write $\mu$, variance $\sigma^2$ and SD $\sigma$ with units.
\item Translate the wording into an event in $X$. For each bound compute $z=(x-\mu)/\sigma$, preserving its sign.
\item Because $\sigma>0$, inequality directions stay the same. Use left-cumulative table areas, subtraction, symmetry or a complement.
\item Interpret the answer in the original context and explain the z-score as a signed number of standard deviations, not a currency amount or probability.
\end{enumerate}
\subsection{Table excerpt and normal notation}
Throughout, $N(\mu,\sigma^2)$ gives \textbf{mean and variance}, not mean and SD. The table gives $\Phi(z)=P(Z\le z)$. Use $\Phi(-z)=1-\Phi(z)$.
\[
\begin{array}{c|rrrrrrrr}
z&0&0.50&0.80&1.00&1.20&1.50&2.00&3.00\\ \hline
\Phi(z)&.5000&.6915&.7881&.8413&.8849&.9332&.9772&.9987
\end{array}
\]
The normal model is an illustrative assumption for the relevant range, not a guarantee about actual future spending or delays.
\section{General Worked Problems}\label{sec:general-worked}
\SectionProblemsTableOfContents{
\SectionProblemLink{prob:sec:general-worked0}{Transfer-processing times}\\
\SectionProblemLink{prob:sec:general-worked1}{A spending band and its complement}\\
\SectionProblemLink{prob:sec:general-worked2}{Comparing relative financial performance}}

\subsection{Transfer-processing times}\label{prob:sec:general-worked0}
Processing time $T\sim N(56,5^2)$ minutes for a bank service. Normality is assumed for this illustrative operational model.
\QuestionsTasks[sec:general-worked]{{6}/{sol:sec:general-worked0}/{Find the probability of a processing time above 63.5 minutes and above 52 minutes. Interpret both z-scores.}}
\subsubsection*{Worked solution}\label{sol:sec:general-worked0}
The mean is $\HLPopulationMean{56}$ minutes, variance is $\HLPopulationVariance{25}$ minutes squared and SD is $\HLPopulationSD{5}$ minutes. Standardise:
\[
z_{63.5}=(63.5-56)/5=1.50,\qquad
z_{52}=(52-56)/5=-0.80.
\]
The first boundary is $1.5$ SD above average; the second is $0.8$ SD below.
\[
P(T>63.5)=1-\Phi(1.50)=1-0.9332=\HLProbability{0.0668},
\]
\[
P(T>52)=1-\Phi(-0.80)=\Phi(0.80)=\HLProbability{0.7881}.
\]
Thus about $6.68\%$ exceed 63.5 minutes and $78.81\%$ exceed 52 minutes under the model. A negative standardised boundary can still have a large upper-tail probability.
\secInicio[prob:sec:general-worked0][sec:general-worked]


\subsection{A spending band and its complement}\label{prob:sec:general-worked1}
Household weekly spending is modeled by $X\sim N(100,400)$ dollars, where 400 is the variance in dollars squared.
\QuestionsTasks[sec:general-worked]{{6}/{sol:sec:general-worked1}/{Find the probability of spending between 80 and 130 dollars and of spending outside that band. Explain why dividing by 400 would be wrong.}}
\subsubsection*{Worked solution}\label{sol:sec:general-worked1}
First recover $\HLPopulationSD{\sigma}=\sqrt{400}=\HLPopulationSD{20}$ dollars. The two standardised bounds are
\[
(80-100)/20=-1,\qquad (130-100)/20=1.5.
\]
Therefore
\[
P(80<X<130)=\Phi(1.5)-\Phi(-1)
=0.9332-(1-0.8413)=\HLProbability{0.7745}.
\]
The probability outside is $1-0.7745=\HLTailArea{0.2255}$. About $77.45\%$ of modeled household-weeks fall in the band. Variance is measured in dollars squared; dividing a dollar difference by it would not produce a unitless standard score.
\secInicio[prob:sec:general-worked1][sec:general-worked]


\subsection{Comparing relative financial performance}\label{prob:sec:general-worked2}
Two annual return models, in percentage-point units, are $R_A\sim N(6,2^2)$ and $R_B\sim N(8,4^2)$. Both record an observed return of 10 percentage points.
\QuestionsTasks[sec:general-worked]{{6}/{sol:sec:general-worked2}/{Compare the observed standard scores and the modeled probabilities of a return above 10. Explain why the same observed return can represent different relative performance.}}
\subsubsection*{Worked solution}\label{sol:sec:general-worked2}
For A, $z_A=(10-6)/2=\HLCriticalValue{2}$; for B, $z_B=(10-8)/4=\HLCriticalValue{0.5}$. The common return is two SD above A's mean but only half an SD above B's mean.
\[
P(R_A>10)=1-0.9772=\HLProbability{0.0228},\qquad
P(R_B>10)=1-0.6915=\HLProbability{0.3085}.
\]
A's observation is more unusual relative to its own model. B's chance of exceeding 10 is higher because its mean is higher and its spread is larger. The z-score compares relative location, not the desirability of every feature of an investment.
\secInicio[prob:sec:general-worked2][sec:general-worked]
\section{Advanced Problems}\label{sec:advanced}
\SectionProblemsTableOfContents{
\SectionProblemLink{prob:sec:advanced0}{A changed mean and a changed spread}}

\subsection{A changed mean and a changed spread}\label{prob:sec:advanced0}
Monthly spending is modeled by $X\sim N(100,20^2)$ dollars. A revised policy model is $Y\sim N(110,10^2)$ dollars. A household wants spending above 130 dollars to be less likely.
\QuestionsTasks[sec:advanced]{{6}/{sol:sec:advanced0}/{Compare the probabilities above 130 and below 90 dollars in both models. Justify how a higher mean can coexist with a smaller upper-tail probability at the chosen boundary.}}
\subsubsection*{Worked solution}\label{sol:sec:advanced0}
For the upper boundary, $z_X=(130-100)/20=1.5$ and $z_Y=(130-110)/10=2$. Thus
\[
P(X>130)=1-0.9332=\HLProbability{0.0668},\qquad
P(Y>130)=1-0.9772=\HLProbability{0.0228}.
\]
For the lower boundary, $z_X=(90-100)/20=-0.5$ and $z_Y=(90-110)/10=-2$:
\[
P(X<90)=1-0.6915=\HLProbability{0.3085},\quad
P(Y<90)=1-0.9772=\HLProbability{0.0228}.
\]
The revised distribution is centred ten dollars higher but its SD is halved. At 130 dollars the boundary is farther above its mean in SD units, reducing the upper tail. Lower-than-90 spending also becomes less likely. The stated upper-tail goal is met under the assumed models, but this does not imply that spending is generally lower.
\secInicio[prob:sec:advanced0][sec:advanced]

\end{document}

